% page 479 du fac-simile (image p-478 du scan)
% bloc 170.2 x 254.9 mm ; marge gauche 31.8 mm
\pgc{10.160mm}{0.000mm}{
\l{\cel{50}471}
\l{}
\l{\sou{pudlar}. (\sou{trans.}) (\sou{tekn.}) Rafinar gisfero per deprenar lua troa}
\l{\cel{3}karbono, transformar lu a fero o stalo, per varmigar lua sur-}
\l{\cel{3}faci oxidigiva. - DEFR.}
\l{}
\l{\sou{pudorar}. (\sou{netran}s.(\sou{precipue pri la homini}). Shamar pri to quo}
\l{\cel{3}koncernas sexuo. - EFIS.}
\l{}
\l{\sou{pudro}. Amilo, reduktita a partikuli extreme tenua, parfumizita,}
\l{\cel{3}uzata por la karnaciono, la hari. - DEFR.}
\l{}
\l{\sou{puero}. Homo, kande lu evas inter 7 e 15 yari. - EFL.}
\l{}
\l{\sou{pufo}. Grosa tabureto cilindra, polsterizita, sen apogilo. - DER.}
\l{}
\l{\sou{pugno}. La manuo kande klozita, do kun la fingri interklemanta.-}
\l{\cel{3}EFIS.}
\l{}
\l{\sou{pulo}. (\sou{Lor ludo per karti, biliardo, en la parii pri kavalkuri)}.}
\l{\cel{3}Pario\cel{2}en qua la riskaji (pekunio) di omna ludanti esos tote}
\l{\cel{3}por la ganonto.- EF.}
\l{}
\l{\sou{pulco}. (\sou{zool.})\cel{2}Mikra insekto, parazito di la homo e di ula ani-}
\l{\cel{3}mali, qua su nutras ek lia sango. - L.\cel{1}\sou{pulex irritans}.-FIS.}
\l{}
\l{\sou{pulchinelo}. I. Rolo komika, reprezentata kun du gibi (un dorse,}
\l{\cel{3}un ventre), nazo hokatra e koloroza, kostumo tre precoza, e mi-}
\l{\cel{3}maskilo nigra. - II. (\sou{metaf.}) Persono groteska, poke valoroza. -}
\l{\cel{3}DEFIRS.}
\l{}
\l{\sou{pulio}. (\sou{tekn.}) Roto, ligna o metala, qua rotacas per axo, e pro-}
\l{\cel{3}vizita ye lua cirklokurvo, ye ranuro en qua pasas kordo di qua}
\l{\cel{3}la extremaji recevas : ica, la tiro-forco, ita la rezisto-forco}
\l{\cel{3}multa-kaze, sur transmiso-pulio, rimeno remplasas la kordo, e}
\l{\cel{3}la ranuri divenas neutila. - EFIS.}
\l{}
\l{\sou{pulikario}. (\sou{bot.)}\cel{3}Genero de kompozaji, tribuo de la radicei,}
\l{\cel{3}- L.\cel{1}\sou{pulicaria}. - L.}
\l{}
\l{\sou{pulmono}. (\sou{anat.)}\cel{2}Singla ek la du viceri sponjatra, qui esas in-}
\l{\cel{3}tersimetre en la kavajo torakala, per qui agesas e la respiro, e}
\l{\cel{3}la produkto di voco. - EFIS.}
\l{}
\l{\sou{pulmonario}. (\sou{bot.)}\cel{2}Nomo diplanti diversa qui aspektas, plu o min}
\l{\cel{3}multe, quale pulmono. - L.\cel{1}\sou{pulmonaria.}}
\l{}
\l{\sou{pulmonito.}\cel{1}(\sou{patol.)}\cel{2}Inflamuro di la parenkimo di la pulmono.- FIS}
\l{}
\l{\sou{pulpo}. I. (\sou{bot.}) Parto quaze-karna di ula frukti, legumi, e c. -}
\l{\cel{3}II. (\sou{farmacio}).Pulpo pastigita por ula preparaji farmaciala.-}
\l{\cel{3}- DEFIS.}
\l{}
\l{pulsar. (trans.)I. Diplasar per preso, stroko, shoko. - II. (\sou{pri}}
\l{\cel{3}kordio).Pulsar sango per stroki iterata, sen-cesa. -DEFIRS.}
\l{}
\l{\sou{pultro.}\cel{2}Ensemblo ek la uceli quin onu nutras ed edukas en farmo-}
\l{\cel{3}domo, ruro-domo, korto por la animali domestika (hani,gansi,}
\l{\cel{3}anadi, e c.) - E.}
}
